\documentclass[11pt,letterpaper]{article}
\usepackage[T1]{fontenc}
\usepackage{lmodern,microtype}
\usepackage[margin=0.78in,headheight=15pt,headsep=15pt]{geometry}
\usepackage{mathtools,amssymb}
\usepackage{xcolor,booktabs,array}
\usepackage[most]{tcolorbox}
\usepackage{enumitem,fancyhdr,needspace,hyperref}

\definecolor{ink}{HTML}{162332}
\definecolor{navy}{HTML}{17324D}
\definecolor{teal}{HTML}{167D82}
\definecolor{mist}{HTML}{EDF5F4}
\definecolor{muted}{HTML}{5D6B78}
\color{ink}
\hypersetup{colorlinks=true,linkcolor=teal,urlcolor=teal,
  pdfauthor={Igor Rivin},
  pdftitle={Number Theory: Week 3 Problem Set},
  pdfsubject={Congruences, modular inverses, Fermat's little theorem, and continued fractions of quadratic surds}}
\setlength{\parindent}{0pt}
\setlength{\parskip}{4pt}
\setlist[enumerate]{label=(\alph*),leftmargin=1.8em,itemsep=3pt,topsep=4pt}
\setlength{\abovedisplayskip}{7pt}
\setlength{\belowdisplayskip}{7pt}
\newcommand{\Z}{\mathbb Z}
\newcommand{\floor}[1]{\left\lfloor #1\right\rfloor}
\newcommand{\per}[1]{\overline{#1}}
\newcommand{\pts}[1]{\hfill\textcolor{muted}{\small[#1 points]}}
\newcommand{\partpts}[1]{\textcolor{muted}{\small(#1)}}
\newcommand{\edition}{\ifsolutions Instructor solutions\else Student handout\fi}
\newcounter{problem}
\newcommand{\problem}[2]{%
  \par\addvspace{11pt}\Needspace{6\baselineskip}%
  \stepcounter{problem}{\bfseries\color{navy}\theproblem.\ #1}\pts{#2}\par\nobreak}
\newcommand{\topic}[1]{\par\addvspace{9pt}\Needspace{10\baselineskip}{\large\bfseries\color{teal}#1}\par\nobreak}
\newcommand{\studentbreak}{\ifsolutions\else\newpage\fi}
\pagestyle{fancy}
\fancyhf{}
\fancyhead[L]{\small\color{muted}Week 3 \textbar\ Problem Set 3}
\fancyhead[R]{\small\color{muted}Patterns, Proofs, and Primes}
\fancyfoot[L]{\footnotesize\color{muted}\edition}
\fancyfoot[R]{\small\color{teal}\thepage}
\renewcommand{\headrulewidth}{0.3pt}
\NewDocumentEnvironment{solution}{+b}{%
  \ifsolutions
  \begin{tcolorbox}[enhanced,breakable,colback=mist!45,colframe=teal!45,
    boxrule=0.5pt,arc=1mm,left=3mm,right=3mm,top=2mm,bottom=2mm,
    before skip=7pt,after skip=7pt,
    title={Solution},title after break={Solution (continued)},
    colbacktitle=mist,coltitle=teal,fonttitle=\small\bfseries]
  \small #1
  \end{tcolorbox}
  \fi
}{}

\begin{document}
{\LARGE\bfseries\color{navy}Problem Set 3}\hfill
{\small\color{teal}\edition}\\[4pt]
{\large Clock arithmetic, Fermat's little theorem, and repeating continued fractions}\\[4pt]
{\small\color{muted}Week 3 \textbar\ 100 points + 10 optional challenge points}

\ifsolutions
\medskip
\textbf{Teaching note.} Estimated working time is 2--3 hours, excluding the
challenge. Problems 1--5 use only congruences, inverses, and Fermat's little
theorem; the Chinese remainder theorem, Euler's theorem, and the notion of
order are not assumed. Problems 6--8 use the integer-state algorithm, Galois's
criterion, and the period-reversal corollary from the second continued-fraction
lecture. All floors are decided by integer bounds on $\sqrt D$; decimal
approximations should not receive proof credit. Subpart scores appear in
parentheses. Avoid repeatedly penalizing a carried arithmetic error.
\else
\medskip
\textbf{Name:}\ \rule{0.46\linewidth}{0.4pt}\hfill
\textbf{Due:}\ \rule{0.20\linewidth}{0.4pt}
\fi

\medskip
\textbf{Instructions.} Work exactly: no decimal approximations. When you need
the floor of a surd, justify it with integer bounds such as $2<\sqrt7<3$.
Justify every proof or counterexample. You may use B\'ezout's identity,
Euclid's lemma in the form ``if $n\mid cm$ and $\gcd(n,c)=1$, then $n\mid m$,''
and the fact from Problem Set~2 that coprime divisors of $c$ have a product
dividing $c$.

\textbf{Facts from lecture.} \emph{Fermat's little theorem:} if $p$ is prime,
then $a^p\equiv a\pmod p$ for every integer $a$, and
$a^{p-1}\equiv1\pmod p$ when $p\nmid a$.
\emph{Integer states:} if $x_n=(P_n+\sqrt D)/Q_n$ with
$Q_n\mid D-P_n^2$ and $a_n=\floor{x_n}$, then
\[
  P_{n+1}=a_nQ_n-P_n,\qquad Q_{n+1}=\frac{D-P_{n+1}^2}{Q_n},
  \qquad x_{n+1}=\frac{P_{n+1}+\sqrt D}{Q_{n+1}}.
\]
\emph{Galois:} a quadratic surd $x$ has a purely periodic continued fraction
exactly when it is \emph{reduced}: $x>1$ and $-1<x'<0$, where $x'$ is the
conjugate. \emph{Reversal:} if $x=[\per{b_0,\ldots,b_{\ell-1}}]$, then
$-1/x'=[\per{b_{\ell-1},\ldots,b_0}]$.

\topic{I. Congruences and inverses}

\problem{Warming up the clock}{10}
\begin{enumerate}
  \item Find the least nonnegative residue of $-100$ modulo $13$. \partpts{2}
  \item Find the last two digits of $99\cdot101\cdot103$ without multiplying
  the numbers out. \partpts{2}
  \item Prove that every square is congruent to $0$ or $1$ modulo $4$.
  Deduce that no integer congruent to $3$ modulo $4$ is a sum of two
  squares. Is $2027$ a sum of two squares? \partpts{6}
\end{enumerate}
\begin{solution}
\begin{enumerate}
  \item $-100=-8\cdot13+4$, so the residue is $4$.
  \item Modulo $100$ the factors are $-1$, $1$, and $3$. Their product is
  $-3\equiv97$, so the last two digits are $97$. (The product is positive,
  so its last two digits are its residue modulo $100$.)
  \item Every integer is congruent to $0,1,2$, or $3$ modulo $4$, and
  congruences may be squared. The squares of these residues are
  $0,1,4,9\equiv0,1,0,1$. Hence a sum of two squares is congruent to
  $0+0$, $0+1$, or $1+1$, that is, to $0$, $1$, or $2$; the residue $3$
  never occurs. Since $2027=4\cdot506+3$, it is not a sum of two squares.
\end{enumerate}
\end{solution}

\problem{Inverses by Euclid}{14}
\begin{enumerate}
  \item Use the extended Euclidean algorithm to find the inverse of $23$
  modulo $89$. Check your answer by multiplication. \partpts{5}
  \item Solve $23x\equiv5\pmod{89}$. \partpts{2}
  \item Let $n\ge2$. Prove that an inverse modulo $n$, when it exists, is
  unique modulo $n$: if $ab\equiv1$ and $ac\equiv1\pmod n$, then
  $b\equiv c\pmod n$. \partpts{3}
  \item List the units modulo $15$, and give the inverse of each. Which units
  are their own inverses? \partpts{4}
\end{enumerate}
\begin{solution}
\begin{enumerate}
  \item The divisions are
  \[
  89=3\cdot23+20,\qquad 23=1\cdot20+3,\qquad 20=6\cdot3+2,\qquad 3=1\cdot2+1.
  \]
  Working backward,
  \[
  1=3-2=7\cdot3-20=7\cdot23-8\cdot20=31\cdot23-8\cdot89.
  \]
  Hence $23^{-1}\equiv31\pmod{89}$. Check: $23\cdot31=713=8\cdot89+1$.
  \item Multiply by $31$: $x\equiv155\equiv66\pmod{89}$.
  Check: $23\cdot66=1518=17\cdot89+5$. The solution is unique modulo $89$
  because multiplication by the unit $23$ can be undone.
  \item Using only associativity and commutativity of multiplication of
  residues,
  \[
  b\equiv b\cdot1\equiv b(ac)\equiv(ab)c\equiv1\cdot c\equiv c\pmod n.
  \]
  \item The units are the residues coprime to $15$:
  $1,2,4,7,8,11,13,14$. The inverse pairs are
  $2\cdot8=16$, $7\cdot13=91=6\cdot15+1$, and the self-inverse units are
  $1,4,11,14$ ($4^2=16$, $11^2=121=8\cdot15+1$, $14\equiv-1$).
  Note that there are four square roots of $1$ here; compare Problem 3(b).
\end{enumerate}
\end{solution}

\problem{Cancellation, square roots of one, and Wilson's theorem}{14}
\begin{enumerate}
  \item Prove the cancellation law: if $ac\equiv bc\pmod n$ and
  $\gcd(c,n)=1$, then $a\equiv b\pmod n$. Give an example showing that the
  conclusion can fail when $\gcd(c,n)>1$. \partpts{5}
  \item Let $p$ be prime. Prove that $a^2\equiv1\pmod p$ forces
  $a\equiv\pm1\pmod p$. Where does your proof break for the modulus $8$?
  \partpts{4}
  \item Pair each of the residues $2,3,\ldots,9$ with its inverse modulo $11$,
  and use the pairing to compute $10!$ modulo $11$ with almost no arithmetic.
  Then prove \emph{Wilson's theorem}: $(p-1)!\equiv-1\pmod p$ for every
  prime $p$. \partpts{5}
\end{enumerate}
\begin{solution}
\begin{enumerate}
  \item The hypothesis says $n\mid c(a-b)$. Since $\gcd(n,c)=1$, Euclid's
  lemma gives $n\mid a-b$. (Alternatively, multiply both sides by an inverse
  of $c$.) \partpts{3} For a failure, $2\cdot1\equiv2\cdot4\pmod6$ but
  $1\not\equiv4\pmod6$. \partpts{2}
  \item The hypothesis says $p\mid a^2-1=(a-1)(a+1)$. A prime dividing a
  product divides a factor, so $a\equiv1$ or $a\equiv-1\pmod p$.
  \partpts{3} Modulo $8$ we have $3^2\equiv1$ although $3\not\equiv\pm1$:
  here $8\mid2\cdot4$ without dividing either factor, so the step
  ``divides a product, hence divides a factor'' requires a prime. \partpts{1}
  \item Modulo $11$: $2\cdot6=12$, $3\cdot4=12$, $5\cdot9=45=4\cdot11+1$,
  and $7\cdot8=56=5\cdot11+1$. Therefore
  \[
  10!\equiv1\cdot(2\cdot6)(3\cdot4)(5\cdot9)(7\cdot8)\cdot10
    \equiv10\equiv-1\pmod{11}.\quad\partpts{2}
  \]
  In general, the case $p=2$ is $1!\equiv-1\pmod2$. Let $p$ be odd. Each
  of $1,\ldots,p-1$ is coprime to $p$, so it has an inverse in the same
  list, unique by Problem 2(c); and the inverse of the inverse of $a$ is $a$.
  By (b), $a$ is its own inverse exactly when $a\equiv\pm1$, that is, for
  $a=1$ and $a=p-1$ only. The remaining residues $2,\ldots,p-2$ therefore
  split into pairs $\{a,a^{-1}\}$ of distinct elements, each pair with
  product $1$. Hence $(p-1)!\equiv1\cdot(p-1)\equiv-1\pmod p$. \partpts{3}
\end{enumerate}
\end{solution}

\studentbreak
\topic{II. Fermat's little theorem}

\problem{Large powers made small}{12}
\begin{enumerate}
  \item Compute $3^{1000}\pmod{13}$. \partpts{3}
  \item Compute $5^{2026}\pmod{23}$. \partpts{3}
  \item Let $p$ be prime and $p\nmid a$. Explain why $a^{p-2}$ is an inverse
  of $a$ modulo $p$. Use this to find the inverse of $7$ modulo $19$ by
  repeated squaring, and check your answer by multiplication. \partpts{6}
\end{enumerate}
\begin{solution}
\begin{enumerate}
  \item By Fermat, $3^{12}\equiv1\pmod{13}$. Since $1000=12\cdot83+4$,
  \[
  3^{1000}\equiv(3^{12})^{83}\cdot3^4\equiv81\equiv3\pmod{13}.
  \]
  (Even faster: $3^3=27\equiv1$ and $1000=3\cdot333+1$.)
  \item By Fermat, $5^{22}\equiv1\pmod{23}$, and $2026=22\cdot92+2$. Hence
  $5^{2026}\equiv5^2=25\equiv2\pmod{23}$.
  \item $a\cdot a^{p-2}=a^{p-1}\equiv1\pmod p$ by Fermat. \partpts{2}
  For $a=7$, $p=19$ we need $7^{17}=7^{16}\cdot7$. Squaring repeatedly,
  \[
  7^2=49\equiv11,\qquad 7^4\equiv121\equiv7,\qquad 7^8\equiv11,\qquad
  7^{16}\equiv7\pmod{19},
  \]
  so $7^{17}\equiv7\cdot7\equiv11$. Check: $7\cdot11=77=4\cdot19+1$.
  \partpts{4}
\end{enumerate}
\end{solution}

\problem{Fermat as a proof tool, and a warning}{14}
\begin{enumerate}
  \item Prove that $30\mid n^5-n$ for every integer $n$. \partpts{5}
  \item Let $p$ be a prime other than $2$ and $5$. Prove that $p$ divides
  infinitely many of the numbers $9,\,99,\,999,\,9999,\ldots$ \partpts{4}
  \item Verify that $2^{10}=3\cdot341+1$, and deduce that
  $2^{340}\equiv1\pmod{341}$. Yet $341=11\cdot31$ is not prime. Does this
  contradict Fermat's theorem? Finally, use Fermat modulo $31$ to show that
  $3^{340}\not\equiv1\pmod{341}$, and explain how this congruence alone
  proves that $341$ is composite. \partpts{5}
\end{enumerate}
\begin{solution}
\begin{enumerate}
  \item Fermat with $p=5$ gives $n^5\equiv n\pmod5$ directly. Modulo $3$,
  Fermat gives $n^3\equiv n$, so $n^5=n^3\cdot n^2\equiv n\cdot n^2=n^3\equiv n$.
  Modulo $2$, $n^2\equiv n$ gives $n^5=(n^2)^2n\equiv n^2\cdot n\equiv n\cdot n
  \equiv n$. Thus $2$, $3$, and $5$ each divide $n^5-n$. Since
  $\gcd(2,3)=1$, their product $6$ divides $n^5-n$; since $\gcd(6,5)=1$,
  so does $30$. (Award full credit for the factorization
  $n^5-n=(n-1)n(n+1)(n^2+1)$ with a complete case analysis modulo $5$.)
  \item The number with $m$ nines is $10^m-1$. Since $p\nmid10$, Fermat
  gives $10^{p-1}\equiv1\pmod p$, hence $10^{k(p-1)}\equiv1$ for every
  $k\ge1$. Thus $p$ divides $10^{k(p-1)}-1$, the number with $k(p-1)$
  nines, for every $k\ge1$; these are infinitely many different numbers.
  \item $3\cdot341=1023=2^{10}-1$. Hence $2^{10}\equiv1$ and
  $2^{340}=(2^{10})^{34}\equiv1\pmod{341}$. \partpts{2} There is no
  contradiction: Fermat's theorem says that primes satisfy the
  congruence; it does not say that only primes do. The \emph{converse} of
  Fermat's theorem is false. \partpts{1} For base $3$: by Fermat,
  $3^{30}\equiv1\pmod{31}$, and $340=30\cdot11+10$, so
  $3^{340}\equiv3^{10}\pmod{31}$. Now $3^3=27\equiv-4$, so
  $3^9\equiv-64\equiv-2$ and $3^{10}\equiv-6\equiv25\pmod{31}$.
  If $3^{340}$ were $1$ modulo $341$, it would also be $1$ modulo the
  divisor $31$; it is not. Finally, if $341$ were prime, Fermat (with
  $341\nmid3$) would force $3^{340}\equiv1\pmod{341}$. The failure of this
  congruence proves compositeness without exhibiting a factor.
  \partpts{2} (In fact $3^{340}\equiv56\pmod{341}$.)
\end{enumerate}
\end{solution}

\studentbreak
\topic{III. Continued fractions of quadratic surds}

\problem{Running the integer-state algorithm}{14}
\begin{enumerate}
  \item Make a table of $n$, $P_n$, $Q_n$, $x_n$, $a_n$ for $x_0=\sqrt7$,
  continuing until a state $(P_n,Q_n)$ repeats. State the continued
  fraction of $\sqrt7$. \partpts{6}
  \item Let $x=(1+\sqrt2)/3$. Explain why $(P_0,Q_0)=(1,3)$ with $D=2$ is
  not a legal starting state. Show that $x=(3+\sqrt{18})/9$ is legal, and
  run the algorithm to find the continued fraction of $x$. \partpts{6}
  \item Which of the complete quotients $x_0,x_1,x_2,\ldots$ in (b) are
  reduced? Compare with where the period begins. \partpts{2}
\end{enumerate}
\begin{solution}
\begin{enumerate}
  \item Every floor follows from $2<\sqrt7<3$.
  \begin{center}
  \begin{tabular}{@{}ccccc@{}}
  \toprule
  $n$ & $P_n$ & $Q_n$ & $x_n$ & $a_n$\\
  \midrule
  $0$ & $0$ & $1$ & $\sqrt7$ & $2$\\
  $1$ & $2$ & $3$ & $(2+\sqrt7)/3$ & $1$\\
  $2$ & $1$ & $2$ & $(1+\sqrt7)/2$ & $1$\\
  $3$ & $1$ & $3$ & $(1+\sqrt7)/3$ & $1$\\
  $4$ & $2$ & $1$ & $2+\sqrt7$ & $4$\\
  $5$ & $2$ & $3$ & $(2+\sqrt7)/3$ & $1$\\
  \bottomrule
  \end{tabular}
  \end{center}
  For instance, $4<2+\sqrt7<5$ gives $4/3<x_1<5/3$, so $a_1=1$; then
  $P_2=1\cdot3-2=1$ and $Q_2=(7-1)/3=2$. Since $x_5=x_1$,
  \[
  \sqrt7=[2;\per{1,1,1,4}].
  \]
  \item A legal state needs $Q_0\mid D-P_0^2$, but $3\nmid2-1=1$; the
  recurrence would produce the non-integer $Q_1$. Multiplying numerator
  and denominator by $3$ gives $x=(3+\sqrt{18})/9$, and
  $9\mid18-9=9$. \partpts{2} Using $4<\sqrt{18}<5$:
  \begin{center}
  \begin{tabular}{@{}ccccc@{}}
  \toprule
  $n$ & $P_n$ & $Q_n$ & $x_n$ & $a_n$\\
  \midrule
  $0$ & $3$ & $9$ & $(3+\sqrt{18})/9$ & $0$\\
  $1$ & $-3$ & $1$ & $-3+\sqrt{18}$ & $1$\\
  $2$ & $4$ & $2$ & $(4+\sqrt{18})/2$ & $4$\\
  $3$ & $4$ & $1$ & $4+\sqrt{18}$ & $8$\\
  $4$ & $4$ & $2$ & $(4+\sqrt{18})/2$ & $4$\\
  \bottomrule
  \end{tabular}
  \end{center}
  Here $7/9<x_0<8/9$, $1<x_1<2$, $4<x_2<9/2$, and $8<x_3<9$. Since
  $x_4=x_2$,
  \[
  \frac{1+\sqrt2}{3}=[0;1,\per{4,8}].\quad\partpts{4}
  \]
  \item $x_0<1$ is not reduced. $x_1>1$, but its conjugate
  $-3-\sqrt{18}$ is below $-1$, so $x_1$ is not reduced. The conjugate of
  $x_2$ is $(4-\sqrt{18})/2$, which lies in $(-1/2,0)$, and the conjugate
  of $x_3$ is $4-\sqrt{18}\in(-1,0)$; both exceed $1$, so both are reduced,
  as is every later $x_n$. The reduced complete quotients are exactly the
  ones inside the period, as Galois's theorem requires: $x_n$ is purely
  periodic precisely for $n\ge2$.
\end{enumerate}
\end{solution}

\problem{From a period back to a surd}{12}
\begin{enumerate}
  \item Let $x=[\per{3,1}]$. Find a quadratic equation with integer
  coefficients satisfied by $x$, and write $x$ in the form
  $(P+\sqrt D)/Q$. Confirm directly that $x$ is reduced. \partpts{5}
  \item Use the reversal corollary to find $t=[\per{1,3}]$ in closed form,
  and then find $y=[2;\per{1,3}]$. \partpts{4}
  \item Your answers should satisfy $y=x-1$. Explain this by looking only
  at the digits. Is $y$ purely periodic? Reconcile your answer with
  Galois's theorem. \partpts{3}
\end{enumerate}
\begin{solution}
\begin{enumerate}
  \item Removing one period leaves $x$ again:
  \[
  x=3+\cfrac1{1+\cfrac1x}=3+\frac{x}{x+1}=\frac{4x+3}{x+1}.
  \]
  (In matrix form, $D(3)D(1)=\begin{psmallmatrix}4&3\\1&1\end{psmallmatrix}$.)
  Hence $x^2-3x-3=0$. Since $x>0$,
  \[
  x=\frac{3+\sqrt{21}}2.
  \]
  From $4<\sqrt{21}<5$ we get $7/2<x<4$, so $x>1$, and
  $x'=(3-\sqrt{21})/2$ lies between $-1$ and $-1/2$. Thus $x$ is reduced.
  \item By the corollary,
  \[
  t=-\frac1{x'}=\frac{2}{\sqrt{21}-3}
   =\frac{2(\sqrt{21}+3)}{12}=\frac{3+\sqrt{21}}6.
  \]
  Then
  \[
  y=2+\frac1t=2+\frac{6}{3+\sqrt{21}}
   =2+\frac{6(\sqrt{21}-3)}{12}=\frac{1+\sqrt{21}}2.
  \]
  \item Writing out the period, $x=[3;1,3,1,3,\ldots]=[3;\per{1,3}]$,
  whereas $y=[2;\per{1,3}]$. The two expansions have the same tail
  $t$ and first digits differing by one: $x=3+1/t$ and $y=2+1/t$.
  \partpts{1} The number $y$ is not purely periodic: the conjugate
  $y'=(1-\sqrt{21})/2$ satisfies $y'<(1-4)/2<-1$, so $y$ is not reduced,
  and Galois's theorem excludes pure periodicity. \partpts{2} (Contrast
  the lecture example $[2;\per{1,2}]=[\per{2,1}]$, where the displayed
  prefix was not a genuine preperiod. Here it is genuine: a purely
  periodic expansion of $y$ with first digit $2$ would have to be
  $[\per{2,1,3,\ldots}]$, and uniqueness of the expansion would force the
  digit $2$ to recur.)
\end{enumerate}
\end{solution}

\problem{Galois's test and two infinite families}{10}
\begin{enumerate}
  \item Without computing any continued fractions, decide which of
  \[
  2+\sqrt7,\qquad \frac{3+\sqrt5}2,\qquad \frac{2+\sqrt7}3,\qquad
  \frac{1+\sqrt7}4
  \]
  have purely periodic continued fractions. Then read off the expansions of
  the purely periodic ones from your table in Problem 6(a). \partpts{4}
  \item Let $n\ge1$. Prove that $\sqrt{n^2+1}=[n;\per{2n}]$. \partpts{3}
  \item Let $n\ge1$. Find and prove a formula for the continued fraction of
  $\sqrt{n^2+2}$. \partpts{3}
\end{enumerate}
\begin{solution}
\begin{enumerate}
  \item Use $2<\sqrt7<3$ and $2<\sqrt5<3$.
  \begin{center}
  \begin{tabular}{@{}llll@{}}
  \toprule
  $x$ & $x>1$? & $x'$ & Reduced?\\
  \midrule
  $2+\sqrt7$ & yes & $2-\sqrt7\in(-1,0)$ & yes\\
  $(3+\sqrt5)/2$ & yes & $(3-\sqrt5)/2>0$ & no\\
  $(2+\sqrt7)/3$ & yes, $x>4/3$ & $(2-\sqrt7)/3\in(-1/3,0)$ & yes\\
  $(1+\sqrt7)/4$ & no, $x<1$ & & no\\
  \bottomrule
  \end{tabular}
  \end{center}
  The reduced ones are the complete quotients $x_4$ and $x_1$ of $\sqrt7$.
  Reading the digits from those rows onward,
  \[
  2+\sqrt7=[\per{4,1,1,1}],\qquad \frac{2+\sqrt7}3=[\per{1,1,1,4}].
  \]
  \item Put $D=n^2+1$. From $n^2<D<(n+1)^2$ we get $n<\sqrt D<n+1$, so
  $a_0=n$ and
  \[
  x_1=\frac1{\sqrt D-n}=\frac{\sqrt D+n}{D-n^2}=\sqrt D+n .
  \]
  Since $2n<x_1<2n+1$, $a_1=2n$ and $x_2=1/(\sqrt D-n)=x_1$. Every later
  digit is therefore $2n$.
  \item The formula is $\sqrt{n^2+2}=[n;\per{n,2n}]$. Put $D=n^2+2$. Since
  $n\ge1$ gives $n^2<D<n^2+2n+1$, again $n<\sqrt D<n+1$ and $a_0=n$.
  The states are
  \[
  (P_1,Q_1)=(n,2),\quad x_1=\frac{n+\sqrt D}2\in\Bigl(n,\,n+\tfrac12\Bigr),
  \quad a_1=n;
  \]
  \[
  (P_2,Q_2)=\Bigl(2n-n,\ \frac{D-n^2}{2}\Bigr)=(n,1),\quad
  x_2=n+\sqrt D\in(2n,2n+1),\quad a_2=2n;
  \]
  and then $(P_3,Q_3)=(2n-n,\,(D-n^2)/1)=(n,2)=(P_1,Q_1)$. The states
  repeat with period two. For $n=1$ this recovers
  $\sqrt3=[1;\per{1,2}]$, and $n=2$ gives $\sqrt6=[2;\per{2,4}]$.
\end{enumerate}
\end{solution}

\par\addvspace{12pt}\Needspace{7\baselineskip}
{\bfseries\color{teal}Optional challenge: the period of $\sqrt D$ is a palindrome}
\pts{10 extra}

Let $D$ be a positive nonsquare integer and $a_0=\floor{\sqrt D}$. Prove that
\[
  \sqrt D=[a_0;\per{a_1,a_2,\ldots,a_{\ell-1},2a_0}]
  \qquad\text{with}\qquad a_i=a_{\ell-i}\quad(1\le i\le\ell-1).
\]
Compare with $\sqrt7$ in Problem 6. \emph{Hint:} apply Galois's theorem and
the reversal corollary to $x=a_0+\sqrt D$. You may use the fact that an
irrational number has only one simple continued fraction expansion.
\begin{solution}
Let $x=a_0+\sqrt D$. Then $x>1$, and $x'=a_0-\sqrt D\in(-1,0)$ because
$a_0<\sqrt D<a_0+1$. By Galois's theorem $x$ is purely periodic, and its
first digit is $\floor x=2a_0$. Let $\ell$ be a period length and write
\[
  x=[\per{2a_0,a_1,\ldots,a_{\ell-1}}].
\]
Subtracting $a_0$ changes only the first digit, so
$\sqrt D=[a_0;a_1,\ldots,a_{\ell-1},2a_0,a_1,\ldots]
=[a_0;\per{a_1,\ldots,a_{\ell-1},2a_0}]$, which is the first claim.

Now compute $-1/x'$ in two ways. Directly,
$-1/x'=1/(\sqrt D-a_0)$, which is the first complete quotient of $x$;
deleting the first digit of $x$ gives
\[
  -\frac1{x'}=[\per{a_1,a_2,\ldots,a_{\ell-1},2a_0}].
\]
By the reversal corollary applied to $x$,
\[
  -\frac1{x'}=[\per{a_{\ell-1},\ldots,a_2,a_1,2a_0}].
\]
By uniqueness of the expansion, the two digit sequences agree term by term:
$a_i=a_{\ell-i}$ for $1\le i\le\ell-1$. For $\sqrt7=[2;\per{1,1,1,4}]$ the
palindrome is $1,1,1$ and the last digit is $4=2\cdot2$.

\textbf{Suggested credit:} $x$ is reduced, hence purely periodic with first
digit $2a_0$ (3); shape of the expansion of $\sqrt D$ (2); two expressions
for $-1/x'$ (3); comparison by uniqueness (2).
\end{solution}

\end{document}
